% Preambula: Beamer + Metropolis, bialy ekran, paleta rozu i fioletu, czcionki >= 10 pt
\documentclass[aspectratio=169,12pt]{beamer}
\usetheme[progressbar=none,numbering=none,background=light]{metropolis}
\usepackage{fontspec}
\usepackage{polyglossia}
\setdefaultlanguage{polish}
\usepackage{xcolor}
\usepackage{tikz}
\usetikzlibrary{arrows.meta,positioning,shapes.misc,fit,calc,backgrounds}
\usepackage{listings}
\usepackage{amsmath}

% --- czcionki (z zapasowymi) ---
\IfFontExistsTF{Inter}{\setsansfont{Inter}}{\IfFontExistsTF{Fira Sans}{\setsansfont{Fira Sans}}{\setsansfont{TeX Gyre Heros}}}
\IfFontExistsTF{Latin Modern Mono}{\setmonofont{Latin Modern Mono}}{\IfFontExistsTF{DejaVu Sans Mono}{\setmonofont{DejaVu Sans Mono}}{}}

% --- paleta ---
\definecolor{plum}{HTML}{7B2D5F}
\definecolor{rose}{HTML}{D94F8A}
\definecolor{violet}{HTML}{6C5CB8}
\definecolor{lav}{HTML}{B7A9E0}
\definecolor{blush}{HTML}{FBE4EE}
\definecolor{lilac}{HTML}{EDE7F6}
\definecolor{ink}{HTML}{3B2A3F}
\definecolor{grey}{HTML}{7C7180}
\definecolor{codebg}{HTML}{FDF3F8}
\setbeamercolor{normal text}{fg=ink,bg=white}
\setbeamercolor{background canvas}{bg=white}
\setbeamercolor{frametitle}{fg=plum,bg=white}
\setbeamercolor{alerted text}{fg=rose}
\setbeamercolor{itemize item}{fg=rose}
\setbeamercolor{itemize subitem}{fg=rose}
\setbeamerfont{itemize/enumerate body}{size=\normalsize}
\setbeamertemplate{itemize item}{\raisebox{0.1ex}{\fontsize{10}{10}\selectfont$\blacksquare$}}
\setbeamersize{text margin left=0.5cm,text margin right=0.5cm}

% --- tytul ramki: wyzej, z cienka linia ---
\setbeamertemplate{frametitle}{%
  \vspace*{-0.12cm}%
  {\color{plum}\bfseries\fontsize{17}{19}\selectfont\insertframetitle\par}%
  \vspace*{-0.2cm}{\color{rose}\rule{\textwidth}{1.1pt}}\par\vspace*{0.05cm}}
\setbeamertemplate{headline}{}
\newcommand{\fsrc}{}
\newcommand{\setsrc}[1]{\gdef\fsrc{#1}}
\setbeamertemplate{footline}{%
  \hspace*{0.5cm}{\color{grey}\fontsize{10}{12}\selectfont\fsrc}\hfill
  {\color{grey}\fontsize{10}{12}\selectfont\insertframenumber\,/\,32}\hspace*{0.5cm}\vspace*{0.18cm}}
\setbeamertemplate{navigation symbols}{}

% --- kod i konsola ---
\lstdefinestyle{pycode}{language=Python,basicstyle=\ttfamily\fontsize{10}{12}\selectfont\color{ink},
  keywordstyle=\bfseries\color{rose},stringstyle=\color{violet},commentstyle=\itshape\color{grey},
  showstringspaces=false,backgroundcolor=\color{codebg},frame=single,rulecolor=\color{blush},
  framesep=3pt,xleftmargin=3pt,xrightmargin=3pt,columns=fixed,keepspaces=true,
  morekeywords={as},emph={self},aboveskip=0pt,belowskip=0pt}
\lstdefinestyle{console}{language={},basicstyle=\ttfamily\fontsize{10}{12}\selectfont\color{plum},
  backgroundcolor=\color{white},frame=l,rulecolor=\color{rose},framerule=2pt,framesep=4pt,
  xleftmargin=5pt,columns=fixed,keepspaces=true,showstringspaces=false,aboveskip=0pt,belowskip=0pt}
\lstdefinestyle{shell}{style=pycode,language=bash,morekeywords={python,pip,source}}
\newcommand{\codeblock}[1]{\lstinputlisting[style=pycode]{code/#1}}
\newcommand{\shellblock}[1]{\lstinputlisting[style=shell]{code/#1}}
\newcommand{\consblock}[1]{\lstinputlisting[style=console]{console/#1}}
\newcommand{\fig}[2]{\includegraphics[#1]{figures/#2}}
\newcommand{\msg}[1]{{\color{rose}\bfseries #1}\par\vspace{4pt}}
\newcommand{\hl}[1]{{\color{rose}\bfseries #1}}

% --- schematy TikZ ---
\tikzset{
  bx/.style={draw=plum,rounded corners=3pt,fill=blush,text=ink,align=center,minimum height=0.95cm,
             inner sep=3pt,font=\fontsize{11}{13}\selectfont,line width=0.8pt},
  bv/.style={bx,draw=violet,fill=lilac},
  bg/.style={bx,draw=grey,fill=white},
  ar/.style={-{Stealth[length=2.6mm]},line width=1.1pt,color=plum},
  lbl/.style={font=\fontsize{10}{12}\selectfont,text=grey,align=center},
}

% --- uklad: rysunek u gory, kod i wynik pod spodem ---
\newcommand{\tslide}[2]{%
  \centering\fig{width=\linewidth}{#1}\par\vspace{5pt}
  \begin{columns}[T,onlytextwidth]
  \begin{column}{0.70\textwidth}\codeblock{#2.py}\end{column}
  \begin{column}{0.28\textwidth}\consblock{#2.txt}\end{column}
  \end{columns}}
